\chapter{Lendo bytecode para entender bugs e decisões}
\label{cap:lendo-bytecode}

\proximasaulas

Você não precisa escrever bytecode. Precisa conseguir \emph{ler} o
suficiente para responder perguntas como: por que este \texttt{==} entre
\texttt{Integer} deu \texttt{false}? Por que o \texttt{switch} lançou
\texttt{NullPointerException}? Mudei uma constante e o outro módulo não
viu, por quê? Cada seção deste capítulo segue o mesmo roteiro: o código,
a listagem real do \texttt{javap} (JDK~25), a leitura das instruções que
importam, o bug ou a decisão que elas explicam e, quando há ligação
documentada, uma ponte com o Quarkus.

O texto avisa quando o bytecode muda entre versões do JDK: a
concatenação de strings via \texttt{invokedynamic} existe desde o JDK~9
e o pattern matching em \texttt{switch} desde o JDK~21.

\section{O básico do \texttt{javap}}
\label{sec:javap-basico}

O \texttt{javap} vem com o JDK e desmonta arquivos \texttt{.class}
(\citetitle{doc:javap}). As opções que usamos:

\begin{description}
\item[\texttt{-c}] imprime as instruções de bytecode de cada método;
\item[\texttt{-p}] mostra também membros privados (sem ela, o método
  sintético de um lambda some);
\item[\texttt{-v}] informação adicional: versão do \texttt{.class},
  constant pool, atributos, tamanho da pilha;
\item[\texttt{-l}] tabelas de linhas e variáveis locais (exige compilar
  com \texttt{-g}).
\end{description}

Uso típico: \texttt{javap -c -p -cp target/classes com.exemplo.MinhaClasse}.

\paragraph{Uma máquina de pilha.}
A JVM não tem registradores visíveis: cada chamada de método ganha um
\emph{frame} com um vetor de variáveis locais e uma pilha de operandos
(JVMS §2.6). As instruções empilham, operam no topo da pilha e
desempilham. Num método de instância, a variável local 0 é
\texttt{this}; os parâmetros vêm em seguida.

\lstinputlisting[language=Java]{generated/source/ConcatenacaoStrings.java}

\inputbytecode{ConcatenacaoStrings}

\paragraph{Lendo a listagem.}
\begin{itemize}
\item O construtor que você não escreveu: \texttt{aload\_0} empilha
  \texttt{this}, \texttt{invokespecial} chama
  \texttt{Object.<init>}, \texttt{return}.
\item O prefixo da instrução diz o tipo: \texttt{a} para referência
  (\texttt{aload}, \texttt{areturn}), \texttt{i} para \texttt{int}
  (\texttt{iload}, \texttt{ireturn}). Em \texttt{saudacao}, que é
  \texttt{static}, \texttt{aload\_0} é \texttt{nome} e \texttt{iload\_1}
  é \texttt{idade}.
\item \texttt{\#7}, \texttt{\#11}: índices no constant pool. O
  comentário à direita é o \texttt{javap} resolvendo o índice para você.
\item \nolinkurl{(Ljava/lang/String;I)Ljava/lang/String;} é um
  \emph{descritor} de método (JVMS §4.3): entre parênteses os parâmetros
  (\texttt{L\ldots;} é uma classe, \texttt{I} é \texttt{int}), depois o
  retorno (\texttt{V} é \texttt{void}, \texttt{Z} é \texttt{boolean}).
\item Em \texttt{main}: \texttt{getstatic} pega \texttt{System.out},
  \texttt{ldc} empilha a string \texttt{"Ana"} do constant pool,
  \texttt{bipush 30} empilha o inteiro, \texttt{invokestatic} chama
  \texttt{saudacao} e \texttt{invokevirtual} chama \texttt{println}.
\end{itemize}

\paragraph{As cinco instruções de chamada.}
\texttt{invokestatic} (método estático), \texttt{invokevirtual} (método
de instância, com despacho pelo tipo real), \texttt{invokeinterface}
(método de interface), \texttt{invokespecial} (construtores e
chamadas via \texttt{super}, entre outros casos) e \texttt{invokedynamic}, que a
seção~\ref{sec:invokedynamic} explica.

\section{Constant pool e \texttt{javap -v}}
\label{sec:constant-pool}

Todo \texttt{.class} tem um constant pool (JVMS §4.4): uma tabela com
os nomes de classes, métodos e campos que o código referencia, as
strings literais e as constantes numéricas. As instruções não contêm
nomes, só índices nessa tabela; na resolução (fase de linking, em geral
de forma preguiçosa, no primeiro uso), a JVM troca os índices por
referências concretas (seção~\ref{sec:class-loading}).

\lstinputlisting[language=Java]{generated/source/PoolConstantes.java}

\inputbytecode{PoolConstantes}

\paragraph{Lendo a listagem.}
\begin{itemize}
\item \texttt{major version: 69}: o \texttt{.class} foi gerado para o
  Java~25. Uma JVM mais antiga recusa esse arquivo com
  \texttt{UnsupportedClassVersionError}; é o erro clássico de compilar
  com um JDK e rodar com outro.
\item As entradas se referenciam: \texttt{\#1} é um \texttt{Methodref}
  que aponta para a classe \texttt{\#2} e para o
  \texttt{NameAndType} \texttt{\#3}, que por sua vez aponta para os
  \texttt{Utf8} com o nome \texttt{<init>} e o descritor \texttt{()V}.
\item O campo \texttt{LIMITE} tem o atributo \texttt{ConstantValue: int
  1000}. E em \texttt{rotulo}, offset 1: \texttt{sipush 1000}. O valor
  foi \emph{copiado} para dentro do método; não há \texttt{getstatic}.
\item \texttt{\#9} é um \texttt{InvokeDynamic} que aponta para o
  bootstrap \texttt{\#0}, listado no fim em \texttt{BootstrapMethods}:
  \texttt{StringConcatFactory.makeConcatWithConstants}, com o argumento
  \texttt{"grande: \textbackslash u0001"}. A próxima seção explica.
\end{itemize}

\paragraph{O bug que isso explica.}
Se \texttt{LIMITE} estivesse numa biblioteca e mudasse para 2000, as
classes clientes já compiladas continuariam usando 1000, porque o valor
está no bytecode delas. A JLS (§13.4.9) diz exatamente isso: mudar o
valor de uma constante \texttt{static} não quebra os binários existentes,
mas eles só veem o novo valor se forem recompilados. Na prática:
\enquote{mudei a constante no módulo comum e o serviço continua com o
valor antigo} é sintoma de build incremental que não recompilou os
dependentes.

\paragraph{A decisão que isso muda.}
\texttt{static final} com valor constante (primitivo ou
\texttt{String}) é parte do contrato público de uma biblioteca, como a
assinatura de um método. Valores que podem mudar entre versões (limites
de negócio, URLs, timeouts) devem vir de configuração ou de um método,
não de uma constante compartilhada entre módulos.

\section{\texttt{invokedynamic}: concatenação e lambdas}
\label{sec:invokedynamic}

As outras instruções de chamada sabem, em tempo de compilação, qual
método chamar. O \texttt{invokedynamic} adia essa decisão: na primeira
execução, a JVM chama um \emph{método de bootstrap}, que devolve o código
a usar naquele ponto; nas execuções seguintes, a mesma instrução reutiliza
o resultado, sem chamar o bootstrap de novo (JVMS §5.4.3.6).

\paragraph{Concatenação de strings (JDK~9+).}
Na listagem de \texttt{ConcatenacaoStrings}
(seção~\ref{sec:javap-basico}), \texttt{"Olá, " + nome + "! Idade: " +
idade} virou \emph{uma} instrução \texttt{invokedynamic}, de nome
\texttt{makeConcatWithConstants} e descritor
\nolinkurl{(Ljava/lang/String;I)Ljava/lang/String;}. Só os
valores variáveis vão na pilha; o texto fixo vai como receita no
argumento do bootstrap, com \texttt{\textbackslash u0001} marcando onde
entra cada valor (veja \texttt{PoolConstantes}, na seção anterior). O
\citetitle{doc:jep280} explica o motivo: até o JDK~8 o \texttt{javac}
gerava uma cadeia de \texttt{StringBuilder.append}, e as otimizações do
JIT para reconhecer essa cadeia eram frágeis. Com \texttt{invokedynamic},
a estratégia fica na biblioteca do JDK e pode melhorar sem recompilar o
seu código.

\emph{Bug ou mito:} o conselho \enquote{troque \texttt{+} por
\texttt{StringBuilder}} vem da era do JDK~8. Numa expressão única, o
\texttt{+} já é uma chamada só. Dentro de um laço que acumula numa
\texttt{String}, cada volta ainda cria uma \texttt{String} nova, e aí o
\texttt{StringBuilder} continua sendo a escolha (o tema do item 63 do
\citetitle{livro:bloch2017}).

\paragraph{Lambdas.}
\lstinputlisting[language=Java]{generated/source/LambdaInvokeDynamic.java}

\inputbytecode{LambdaInvokeDynamic}

\begin{itemize}
\item O corpo do lambda virou um método privado e sintético,
  \texttt{lambda\$dobro\$0}. Sem \texttt{-p}, o \texttt{javap} não o
  mostra.
\item Em \texttt{dobro}, uma única instrução: \texttt{invokedynamic}
  \nolinkurl{apply:()Ljava/util/function/Function;}. Com \texttt{-v}, o bootstrap
  aparece como \texttt{LambdaMetafactory.metafactory}, que recebe uma
  referência a \texttt{lambda\$dobro\$0}.
\item Não existe \texttt{LambdaInvokeDynamic\$1.class} no JAR, como
  existiria com uma classe anônima. Segundo o \citetitle{doc:jep371}
  (JDK~15), a \texttt{LambdaMetafactory} define em tempo de execução uma
  \emph{hidden class} que implementa a interface funcional.
\item Dentro de \texttt{lambda\$dobro\$0}: \texttt{intValue}, \texttt{imul},
  \texttt{Integer.valueOf}. Um \texttt{Function<Integer,\allowbreak Integer>}
  desembrulha e reembrulha a cada chamada; a seção seguinte mostra o
  custo e o risco disso.
\end{itemize}

\paragraph{Ponte com Quarkus.}
Executáveis nativos do Quarkus são gerados pelo GraalVM Native Image,
que trabalha com a hipótese de mundo fechado: todo o código é conhecido
em tempo de build. O \citetitle{doc:graalvm-compat} diz que
\texttt{invokedynamic} e method handles em geral são incompatíveis com
essa hipótese, porque podem introduzir chamadas em tempo de execução, mas
que os usos gerados pelo \texttt{javac} para lambdas e concatenação de
strings são suportados, porque não mudam o método chamado. Ou seja:
lambdas e \texttt{+} estão liberados no nativo; geração dinâmica de
código com method handles, não. A Parte~III detalha.

\section{Autoboxing}
\label{sec:autoboxing}

\lstinputlisting[language=Java]{generated/source/Autoboxing.java}

\inputbytecode{Autoboxing}

\paragraph{Lendo a listagem.}
\begin{itemize}
\item Em \texttt{main}, cada \texttt{Integer a = 127} virou
  \texttt{bipush 127} seguido de \texttt{invokestatic Integer.valueOf}.
  O autoboxing é uma chamada de método que o compilador escreveu por
  você.
\item Em \texttt{iguais}, \texttt{a == b} virou \texttt{if\_acmpne}:
  comparação de \emph{referências} (o prefixo \texttt{a}), não de
  valores. Para comparar valores, o compilador teria gerado
  \texttt{intValue} e \texttt{if\_icmpne}.
\end{itemize}

\paragraph{O bug que isso explica.}
O programa imprime \texttt{true} e depois \texttt{false}. A JLS (§5.1.7)
garante que o boxing de uma expressão constante entre $-128$ e $127$
devolve sempre o mesmo objeto; fora dessa faixa, não há garantia, e na
prática 128 gera dois objetos diferentes. O teste passa com valores
pequenos e falha em produção com IDs maiores. Compare wrappers com
\texttt{equals} (ou \texttt{Objects.equals}), nunca com \texttt{==}.

O caminho inverso tem outro risco. Em \texttt{lambda\$dobro\$0}
(seção~\ref{sec:invokedynamic}), \texttt{x * 2} começa com
\texttt{invokevirtual Integer.intValue}. Se \texttt{x} for \texttt{null},
essa chamada lança \texttt{NullPointerException} numa linha que não tem
nenhuma chamada de método visível no código-fonte.

\paragraph{A decisão que isso muda.}
Os itens 6 e 61 do \citetitle{livro:bloch2017} (\enquote{Avoid creating
unnecessary objects} e \enquote{Prefer primitive types to boxed
primitives}) ganham uma explicação concreta: cada \texttt{valueOf} pode
alocar. Em código quente, prefira \texttt{int}, \texttt{IntStream} e
coleções de primitivos; deixe wrappers para onde \texttt{null} tem
significado (campo opcional de um JSON, coluna anulável).

\section{\texttt{switch} sobre \texttt{String}}
\label{sec:switch}

\lstinputlisting[language=Java]{generated/source/SwitchStrings.java}

\inputbytecode{SwitchStrings}

\paragraph{Lendo a listagem.}
\begin{itemize}
\item Offset 5: \texttt{String.hashCode()}. O \texttt{switch} sobre
  \texttt{String} começa calculando o hash.
\item Offset 8: \texttt{lookupswitch} sobre os hashes de
  \texttt{"vermelho"} e \texttt{"verde"}. A JVMS (§3.10) explica a
  escolha: \texttt{tableswitch} quando os casos são densos (um índice
  numa tabela), \texttt{lookupswitch} quando são esparsos (busca numa
  lista ordenada), como valores de hash.
\item Offsets 36--60: hash igual não basta, porque duas strings
  diferentes podem ter o mesmo hash. Cada ramo confirma com
  \texttt{equals} e grava um índice (0, 1) na variável local 2.
\item Offset 62: um \emph{segundo} \texttt{lookupswitch}, agora sobre o
  índice, leva ao código de cada \texttt{case}.
\end{itemize}

\paragraph{O bug que isso explica.}
Se \texttt{cor} for \texttt{null}, o offset 5 chama \texttt{hashCode()}
numa referência nula e lança \texttt{NullPointerException}, antes de
qualquer \texttt{case} ou \texttt{default}. A JLS (§14.11.3) define esse
comportamento; o \texttt{default} não pega \texttt{null}. A partir do
JDK~21 dá para escrever \texttt{case null} (\citetitle{doc:jep441}); no
\texttt{switch} com padrões, o \texttt{javac} chama
\texttt{Objects.requireNonNull} no seletor quando não há
\texttt{case null}, como vimos na seção~\ref{sec:sealed-pattern}.

\paragraph{A decisão que isso muda.}
\texttt{switch} sobre \texttt{String} vindo de fora (parâmetro HTTP,
campo de JSON) precisa tratar \texttt{null} antes, ou usar
\texttt{case null}. Se os valores formam um conjunto fechado, um
\texttt{enum} é melhor: erro de digitação vira erro de compilação, e o
\texttt{switch} sobre \texttt{enum} não depende de \texttt{hashCode}.

\section{\texttt{try}-with-resources}
\label{sec:twr}

\lstinputlisting[language=Java]{generated/source/TryWithResources.java}

\inputbytecode{TryWithResources}

\paragraph{Lendo a listagem.}
\begin{itemize}
\item Offsets 26--30: caminho normal, \texttt{close()} (27) e salto para
  o fim (\texttt{goto 49}, offset 30).
\item A \texttt{Exception table} diz: exceção nos offsets 8 a 25 (o
  \texttt{to}~26 é exclusivo, e o \texttt{close()} em 27 fica fora) vai
  para 33. Ali a exceção principal é guardada (\texttt{astore\_1}) e
  \texttt{close()} é chamado de novo.
\item Se esse \texttt{close()} também lança, a segunda linha da tabela
  leva ao offset 41, que chama
  \texttt{Throwable.addSuppressed}: a exceção do \texttt{close()} é
  pendurada na principal, que é relançada (\texttt{athrow}, offset 48).
\end{itemize}

É a tradução que a JLS (§14.20.3.1) especifica, com a variável
\texttt{\#primaryExc} e \texttt{addSuppressed}. Compare com
\texttt{ThreadsVirtuais} (seção~\ref{sec:virtual-threads}): lá o recurso
vem de uma chamada de método e o \texttt{javac} gerou \texttt{ifnull}
antes de cada \texttt{close()}, como a JLS pede; aqui o recurso vem de
um \texttt{new}, que nunca é \texttt{null}, e o teste não aparece.

\paragraph{O bug que isso explica.}
Com \texttt{try}/\texttt{finally} escrito à mão, uma exceção no
\texttt{close()} substitui a exceção original, e o log mostra
\enquote{erro ao fechar conexão} em vez da causa real. O
\texttt{try}-with-resources preserva a causa e anexa a outra como
\emph{suppressed}: procure \texttt{Suppressed:} no stack trace. É o
argumento do item 9 do \citetitle{livro:bloch2017} (\enquote{Prefer
try-with-resources to try-finally}).

\section{Sincronização}
\label{sec:sincronizacao}

\lstinputlisting[language=Java]{generated/source/Sincronizacao.java}

\inputbytecode{Sincronizacao}

\paragraph{Lendo a listagem.}
\begin{itemize}
\item Offsets 0--3: \texttt{this} é duplicado, guardado na variável
  local 1 e usado por \texttt{monitorenter}.
\item Offsets 4--11: \texttt{contador++}, as três etapas discutidas na
  seção~\ref{sec:modelo-memoria}.
\item Offsets 14--16: \texttt{monitorexit} no caminho normal.
\item Offsets 19--23: o tratador de exceção, de tipo \texttt{any} na
  tabela, libera o monitor e relança. A segunda linha da tabela cobre o
  próprio tratador.
\end{itemize}

A JVMS (§3.14) diz que o compilador garante um \texttt{monitorexit} para
cada \texttt{monitorenter}, termine o método normalmente ou com exceção,
gerando tratadores que capturam qualquer exceção. Ela diz também que um
\emph{método} \texttt{synchronized} normalmente não usa essas instruções:
é marcado com a flag \texttt{ACC\_SYNCHRONIZED}, e a própria chamada
adquire o monitor. Por isso, num método \texttt{synchronized}, o
\texttt{javap -c} não mostra \texttt{monitorenter}; use \texttt{-v} e
procure a flag.

\paragraph{A decisão que isso muda.}
\texttt{synchronized(this)} usa como lock um objeto que qualquer código
com referência a ele também pode usar para sincronizar. Um
\texttt{private final Object lock = new Object()} deixa o lock sob seu
controle. Em JDK~21, \texttt{monitorenter} era também o ponto em que uma
virtual thread ficava presa à carrier ao bloquear; desde o JDK~24
(\citetitle{doc:jep491}) não é mais (seção~\ref{sec:virtual-threads}).

\paragraph{Ponte com Quarkus.}
O guia \citetitle{doc:quarkus-virtual-threads} registra a mesma mudança:
em Java~24 ou mais novo, blocos e métodos \texttt{synchronized} não
causam mais pinning de virtual threads. Antes
de trocar \texttt{synchronized} por \texttt{ReentrantLock} num serviço
Quarkus, confira a versão do JDK em produção.

\section{Ferramentas para olhar a JVM}
\label{sec:ferramentas}

O \texttt{javap} mostra o que o \texttt{javac} gerou. Para ver o que a
JVM faz com isso em execução, estas são as ferramentas, todas com
documentação oficial:

\small
\begin{longtable}{@{}p{3.2cm} p{9cm}@{}}
\toprule
Ferramenta & Para que serve \\
\midrule \endhead
\texttt{javap} & Desmonta \texttt{.class}: \texttt{-c}, \texttt{-p},
  \texttt{-v} (\citetitle{doc:javap}) \\
\texttt{-XX:+\allowbreak PrintCompilation} & Uma linha a cada método compilado pelo
  JIT (\citetitle{doc:java-cmd}) \\
\texttt{-XX:+\allowbreak LogCompilation} & Log detalhado de compilação em
  \texttt{hotspot\_pid<pid>.log} (o manual ainda diz \texttt{hotspot.log});
  exige \texttt{-XX:+UnlockDiagnosticVMOptions} \\
JITWatch & Analisador e visualizador do log de compilação da HotSpot:
  inlining, métodos quentes, bytecode e assembly
  (\citetitle{doc:jitwatch}) \\
JFR & Flight Recorder: grava eventos da aplicação, da JVM e do sistema
  operacional num arquivo, com meta de no máximo 1\% de overhead
  (\citetitle{doc:jep328}, JDK~11). Inicie com
  \texttt{-XX:StartFlightRecording} ou \texttt{jcmd <pid> JFR.start} \\
\texttt{jfr} & Lê gravações: \texttt{print}, \texttt{summary},
  \texttt{view} (\citetitle{doc:jfr-tool}) \\
\texttt{jcmd} & Comandos para uma JVM rodando: \texttt{JFR.start},
  \texttt{Thread.dump\_to\_file}, \texttt{GC.heap\_info},
  \texttt{VM.flags}, \texttt{Compiler.queue} (\citetitle{doc:jcmd}) \\
JMC & JDK Mission Control, projeto do OpenJDK com ferramentas gráficas
  que abrem gravações do JFR (\citetitle{doc:jmc}) \\
async-profiler & Profiler por amostragem de CPU e alocação, que segundo
  a documentação não sofre de \emph{safepoint bias}; gera flame graphs
  (\citetitle{doc:async-profiler}) \\
JMH & Harness para microbenchmarks na JVM (\citetitle{doc:jmh}) \\
\bottomrule
\end{longtable}
\normalsize

\paragraph{Por onde começar.}
Para \enquote{está lento em produção}, comece pelo JFR: o overhead é
baixo, ele já vem no JDK e cobre GC, alocação, locks e CPU. Para
\enquote{este trecho ficou mais rápido?}, use JMH, nunca um laço com
\texttt{nanoTime}. Para \enquote{por que o JIT não otimizou este
método?}, \texttt{PrintCompilation} e depois JITWatch. Para \enquote{o
que exatamente o \texttt{javac} gerou?}, \texttt{javap}.

\section{Onde aprofundar}

\subsection*{Especificações}
A \citetitle{doc:jvms25}: capítulo~3 (como o \texttt{javac} compila
\texttt{switch}, \texttt{synchronized} e exceções), capítulo~4 (formato
do \texttt{.class} e constant pool) e capítulo~6 (cada instrução). A
\citetitle{doc:jls25}: §5.1.7 (boxing), §13.4.9 (constantes e
compatibilidade binária), §14.11 (\texttt{switch}) e §14.20.3
(\texttt{try}-with-resources).

\subsection*{Livros}
\small
\begin{longtable}{@{}p{1.8cm} p{4.4cm} p{2.6cm} p{4cm}@{}}
\toprule
Tema & Livro & Autor(es), ano & Por que ler \\
\midrule \endhead
Java & \citetitle{livro:bloch2017}, 3ª ed. \pt{} \emph{Java Efetivo}
  & \citeauthor{livro:bloch2017}, \citeyear{livro:bloch2017}
  & Itens 6, 9, 61, 63 e 78, citados neste capítulo \\
Desempenho & \citetitle{livro:evans2018} & \citeauthor{livro:evans2018},
  \citeyear{livro:evans2018} & JIT e ferramentas de diagnóstico \\
\bottomrule
\end{longtable}
\normalsize

\subsection*{Vídeos}
\small
\begin{longtable}{@{}p{1.8cm} p{7.2cm} p{3.4cm}@{}}
\toprule
Tema & Vídeo & Quem \\
\midrule \endhead
Lambdas & \citetitle{video:goetz-lambda-hood} (Philly ETE 2014, InfoQ)
  & \citeauthor{video:goetz-lambda-hood}, arquiteto da linguagem Java \\
\bottomrule
\end{longtable}
\normalsize

\subsection*{Ferramentas e documentação}
\citetitle{doc:javap}, \citetitle{doc:java-cmd}, \citetitle{doc:jcmd},
\citetitle{doc:jfr-tool}, \citetitle{doc:jmc}, \citetitle{doc:jitwatch},
\citetitle{doc:async-profiler} e \citetitle{doc:jmh}. Para as pontes:
\citetitle{doc:graalvm-compat} e \citetitle{doc:quarkus-native-tips}.

\section{Desafios práticos}

\desafio{Rode \texttt{javap -c -p} numa classe sua que use lambda e ache
o \texttt{invokedynamic} e o método \texttt{lambda\$\ldots}. Troque o
lambda por uma referência de método (\texttt{Integer::sum}, por exemplo):
o método sintético continua existindo?}

\desafio{Compile o mesmo \texttt{"a" + x} com \texttt{-{}-release 8} (o JDK~25 avisa que a opção está obsoleta) e com
\texttt{-{}-release 25} e compare as listagens. Onde foi parar o
\texttt{StringBuilder}?}

\desafio{Crie duas classes em projetos separados: uma com
\texttt{public static final int LIMITE = 10}, outra que imprime
\texttt{LIMITE}. Mude o valor, recompile só a primeira e rode a segunda.
Depois troque a constante por um método \texttt{limite()} e repita.}

\desafio{Escreva um método \texttt{synchronized} e um bloco
\texttt{synchronized(this)} e compare \texttt{javap -c} e
\texttt{javap -v} dos dois. Onde está o lock em cada caso?}

\desafio{Grave 60 segundos de um serviço seu com
\texttt{jcmd <pid> JFR.start duration=60s filename=app.jfr} e rode
\texttt{jfr summary app.jfr}. Quais eventos aparecem mais? Abra o
arquivo no JMC e ache o método que mais aloca.}

\desafio{Procure no seu projeto um \texttt{==} entre \texttt{Integer},
\texttt{Long} ou outro wrapper. Escreva um teste que falhe por causa dele
antes de corrigir.}

